% Generated by code/interviews-1/03-speed-drills/python/gen_drill.py. Do not edit. French copy by draft/pilot-quant-fr/ivgen_fr.py.
\ivfamily{Trois chiffres fois deux chiffres}{3, 5, 6, 10, 15, 21, 31}{Tester d'abord le chiffre des unités, puis l'ordre de grandeur.}
\ivfamily{Carrés et quasi-carrés}{9, 22, 24, 32, 34}{Vérifier le chiffre des unités, puis arrondir. $(100 + d)(100 - d)$ vaut $10000 - d^2$.}
\ivfamily{Division longue}{7, 13, 16, 17, 29, 30, 35}{Vérifier en multipliant~: le chiffre des unités du quotient multiplié par celui du diviseur doit donner le chiffre des unités du nombre.}
\ivfamily{Pourcentages}{2, 14, 18, 23, 27, 36}{Construire le pourcentage à partir de 10~\%, 5~\% et 1~\%, et s'arrêter dès qu'il ne reste qu'une seule option.}
\ivfamily{Produits de décimaux}{1, 4, 8, 19, 25, 28}{Compter les décimales, puis tester le chiffre des unités du produit.}
\ivfamily{Fractions d'un nombre}{11, 12, 20, 26, 33}{Diviser d'abord par le dénominateur, puis multiplier par le numérateur.}
